\chapter{Task 5 --- Independent Testing and Reproducibility}
\label{ch:task5}

%-------------------------------------------------------

Task 5 asks for a single evaluation of the selected configuration on the
reserved MNIST test set, reporting accuracy, macro-precision,
macro-recall and macro-F1, together with the random seed, the library
versions, the execution environment and instructions to reproduce the
experiments. These are covered in Sections \ref{sec:task5},
\ref{sec:environment} and \ref{sec:repro}.

%%%%%%%%%%%%%%%%%%%%%%%%%%%%%%%%%%%%%%%%%%%%%%%%%%%%%
\section{Test-Set Evaluation}
\label{sec:task5}
%%%%%%%%%%%%%%%%%%%%%%%%%%%%%%%%%%%%%%%%%%%%%%%%%%%%%

Only after the configuration was selected did \texttt{main.py} load the
10,000 test images, evaluating the saved epoch-10 weights of the baseline
once. Listing \ref{lst:test-output} shows the program output and Table
\ref{table:test} the metrics.

\begin{lstlisting}[language={}, numbers=none, caption={Program output: model selection and the single test evaluation.}, label={lst:test-output}]
=== Validation summary after epoch 10
lr_0.01   val acc 0.9868  val F1 0.9867  val loss 0.0473  best val acc 0.9877 (epoch 7)
lr_0.001  val acc 0.9792  val F1 0.9790  val loss 0.0670  best val acc 0.9792 (epoch 10)
lr_0.1    val acc 0.1133  val F1 0.0203  val loss 2.3030  best val acc 0.9338 (epoch 1)
selected configuration: lr_0.01

=== TEST (lr_0.01, evaluated once)
test_loss        0.0402
test_acc         0.9884
macro_precision  0.9882
macro_recall     0.9885
macro_f1         0.9883
\end{lstlisting}

\begin{table}[h]
    \centering
    \caption[Test-set results.]%
    {Test-set results of the selected model (SGD, $\eta = 0.01$, batch
    size 64, weights after epoch 10, seed 42).}

    \begin{tabular}{P{2.4in}|c}
        \specialrule{1.5pt}{0pt}{0pt}
        Metric & Test set (10,000 images) \\
        \hline\hline
        Accuracy & 98.84\% \\
        \specialrule{0.5pt}{0pt}{0pt}
        Macro-precision & 0.9882 \\
        \specialrule{0.5pt}{0pt}{0pt}
        Macro-recall & 0.9885 \\
        \specialrule{0.5pt}{0pt}{0pt}
        Macro-F1 & 0.9883 \\
        \specialrule{0.5pt}{0pt}{0pt}
        Cross-entropy loss & 0.0402 \\
        \specialrule{1.5pt}{0pt}{0pt}
    \end{tabular}
    \label{table:test}
\end{table}

\textbf{Observation.} The model classifies 9,884 of the 10,000 test
images correctly, so it makes 116 errors. The test accuracy (98.84\%)
is close to the validation accuracy (98.68\%), and slightly higher. The
validation set was therefore a good estimate of performance on unseen
data, and selecting on it did not overfit to the validation images.
Macro-precision, macro-recall and macro-F1 are all within 0.0003 of the
accuracy, so no digit is much worse than the others.


%%%%%%%%%%%%%%%%%%%%%%%%%%%%%%%%%%%%%%%%%%%%%%%%%%%%%
\section{Execution Environment and Library Versions}
\label{sec:environment}
%%%%%%%%%%%%%%%%%%%%%%%%%%%%%%%%%%%%%%%%%%%%%%%%%%%%%

\begin{table}[h]
    \centering
    \caption[Software and hardware environment.]%
    {Software and hardware environment used for every experiment.}

    \begin{tabular}{P{1.6in}|P{1.4in}|P{2.4in}}
        \specialrule{1.5pt}{0pt}{0pt}
        Component & Version & Purpose \\
        \hline\hline
        Python & 3.13.5 & Programming language \\
        \specialrule{0.5pt}{0pt}{0pt}
        PyTorch & 2.14.1 (CPU build) & Tensors, layers, autograd, optimizer \citep{paszke2019} \\
        \specialrule{0.5pt}{0pt}{0pt}
        torchvision & 0.29.1 & MNIST download and transforms \\
        \specialrule{0.5pt}{0pt}{0pt}
        scikit-learn & 1.9.1 & Precision, recall and F1 \citep{pedregosa2011} \\
        \specialrule{0.5pt}{0pt}{0pt}
        NumPy / Matplotlib & 2.5.4 / 3.11.2 & Arrays and figures \citep{hunter2007} \\
        \specialrule{0.5pt}{0pt}{0pt}
        Operating system & Windows 11 Home & 64-bit \\
        \specialrule{0.5pt}{0pt}{0pt}
        Hardware & Intel Core i7-13700F, 16~GB RAM & All training on the CPU (16 threads); no GPU used \\
        \specialrule{1.5pt}{0pt}{0pt}
    \end{tabular}
    \label{table:environment}
\end{table}

The computer has an NVIDIA RTX 4060 Ti, but I installed the CPU build of
PyTorch on purpose. The network is small enough to train on the CPU in a
few minutes, and on the CPU \texttt{torch.use\_deterministic\_algorithms(True)}
makes every run give the same numbers, which is harder to guarantee with
GPU kernels.


%%%%%%%%%%%%%%%%%%%%%%%%%%%%%%%%%%%%%%%%%%%%%%%%%%%%%
\section{Instructions to Reproduce the Experiments}
\label{sec:repro}
%%%%%%%%%%%%%%%%%%%%%%%%%%%%%%%%%%%%%%%%%%%%%%%%%%%%%

All results in this report were produced by the code in Appendix
\ref{app:source} on the machine in Table \ref{table:environment}. To
reproduce them:

\begin{enumerate}
  \item Install Python 3.13 and create a virtual environment in the
        project folder:
\begin{lstlisting}[language={}, numbers=none]
python -m venv env
env\Scripts\activate          (Windows)
source env/bin/activate       (Linux / macOS)
\end{lstlisting}
  \item Install the pinned dependencies (CPU build of PyTorch):
\begin{lstlisting}[language={}, numbers=none]
pip install -r requirements.txt
\end{lstlisting}
  \item Run all experiments. MNIST is downloaded to \texttt{./data} on
        the first run (internet needed once):
\begin{lstlisting}[language={}, numbers=none]
python main.py
\end{lstlisting}
        This trains the three configurations (about 10 minutes in total on
        the CPU above), selects one using the validation set, evaluates
        it once on the test set, and writes:
        \begin{itemize}
          \item \texttt{results/output\_log.txt}: the full program output;
          \item \texttt{results/<config>/history.json} and
                \texttt{model.pt}: metrics and final weights per
                configuration;
          \item \texttt{results/test\_metrics.json}: the test metrics;
          \item \texttt{report/figures/*.png}: every figure in this report.
        \end{itemize}
  \item To redo only the selection, the test evaluation and the figures
        from saved results, run \texttt{python main.py --skip-train}.
\end{enumerate}

\textbf{Random seed.} Seed 42 is used for the train/validation split,
the initial weights and the shuffling order (Section \ref{sec:seed}).
Training runs on the CPU with
\texttt{torch.use\_deterministic\_algorithms(True)} and
\texttt{num\_workers=0}, so running \texttt{main.py} again on the same
software versions gives the same numbers. I checked this by training all
three configurations a second time: every epoch-level loss, accuracy and
F1-score, and the final weights, matched the first run exactly.
Different PyTorch versions or a GPU may change the last digits because
of floating-point differences.
